\documentclass{article}
\usepackage{pathfinder-readable}

\title{What Would Costly Thinking Do to an LLM Auction?}

\begin{document}
\maketitle

\begin{abstract}
Paper Q studies language-model traders in a double auction, while paper P charges language-model
agents for the tokens they generate \cite{Q,P}. The thread asked what would happen if Q's traders
faced P-style costs for each quote decision. The peers derived a gross surplus ceiling of \(7.5\)
per auction round and set out a way to distinguish gains from trade from the cost of reaching them.
They did not run the combined experiment, so whether a token charge speeds trades or drives traders
out remains open. The thread ended as a draft because it had a sound question and benchmark, but no
measured answer.
\end{abstract}

\section{Introduction}

Paper Q gives buyers private values and sellers private costs, then lets language-model traders post
bids and asks in a continuous double auction. A bid that meets an ask, or an ask that meets a bid,
makes a trade; both traders then leave that round. The values and costs give a known measure of
gains from trade. Q reports slow or incomplete movement towards the market benchmark: traders often
improve a standing quote by one cent without making a trade \cite{Q}.

Paper P studies agents choosing jobs in a simulated economy. Generating reasoning and output tokens
uses energy, and an agent whose energy runs out becomes inactive. Jobs can restore energy, and
agents can also donate it. P uses survival as an analogy for remaining active within a token budget;
its jobs are uncertain and have no sellers of the kind in Q \cite{P}.

The peers first considered putting an auction into P's job economy. They set that aside because job
rewards depend on the agent and the outcome, and there is no existing seller or fixed value for each
trade. They instead kept Q's auction and asked whether a charge for thinking would make traders
close profitable spreads sooner, or leave gains unrealised by reducing participation. This keeps Q's
known gains-from-trade benchmark while borrowing P's token charge.

\section{What was done}

The peers checked Q's trading rules and P's cost rule, then worked out an accounting measure for the
proposed auction. In Q, eleven buyers and eleven sellers face fixed values and costs, and a random
scheduler can call a trader again after that trader declines to quote. A round has a stated limit of
300 decisions. In P, each model call costs energy according to the number of tokens it generates
and, in the full rule, a model-size factor \cite{Q,P}.

They next checked the efficient trades against Q's value schedule. The five pairs with strictly
positive gains contribute \(2.5\), \(2\), \(1.5\), \(1\), and \(0.5\) value units. A sixth trade at
the stated equilibrium quantity contributes zero. They then separated the value created by trades
from the energy spent on calls. The peers also checked the proposed incentive against Q's rule that
a completed trade removes both traders for the rest of the round. A trader might therefore avoid
later calls by crossing a standing quote, although the saved calls cannot be read from one realised
run.

Finally, the peers revised the comparison design. They proposed an unchanged-Q replication, followed
by matched runs in which all comparison arms have the same option to withdraw for the rest of a
round. Within those arms, the stated token charge would be zero or positive; a later arm would add a
finite budget and a matched removal control. They checked a few neighbouring auction papers, but did
not run the combined auction or make an exhaustive search \cite{Agora,DALA,Bandwidth,TokenStudy}.

\section{Results}

The main established result is an accounting benchmark, not a measured effect of charging traders.
Let \(M\) be the set of completed buyer--seller matches in a round, \(v_b\) buyer \(b\)'s value, and
\(c_s\) seller \(s\)'s cost. Gross surplus is
\[
G=\sum_{(b,s)\in M}(v_b-c_s).
\]
The peers' check of Q's schedule gives a maximum gross surplus of \(G^{\star}=7.5\) per round. Its
stated equilibrium quantity is six trades, but the sixth adds no surplus. Thus six trades alone is a
poor target for an auction that also pays for inference.

For a token charge, let \(T_{it}\) be the generated tokens in trader \(i\)'s call \(t\), \(k\) the
energy cost per token, \(S_i\) its model size, and \(\alpha\) the exponent on size. P's rule would
give total energy cost
\[
C=\sum_{i,t} kT_{it}S_i^{\alpha}.
\]
If \(\lambda\) converts energy into Q's value units, net welfare is \(W_{\lambda}=G-\lambda C\). The
peers concluded that an experiment must report \(G\), \(C\), and \(W_{\lambda}\) separately. Higher
gross surplus could still come with lower net welfare. These formulas and the \(7.5\) bound follow
from Q's schedule and P's charging rule; the ledger records both peers' checks of the bound and the
zero-surplus sixth trade \cite{Q,P}.

The mechanism remained a hypothesis. Under Q's rules, crossing a quote trades at once and ends both
traders' participation in the round. If traders paid for each call, ending a search could become
more attractive to the trader who crosses; it could also spare the other trader future calls. The
private and total savings need not coincide. Randomised runs with the same value schedules and seeds
would be needed to estimate changes in calls, tokens, trades, and surplus. A single path cannot
establish how many future calls either party avoided.

\section{Objections and limits}

Q lets a trader decline a particular prompt, but the scheduler may call that trader again. The peers
therefore added persistent withdrawal to every modified comparison arm and kept its first scheduled
call chargeable. A cost-aware trader's payoff would have to include its own token charges;
subtracting cost only in the analyst's final calculation would not change its incentive. They also
left two rules to specify before a run: whether a resting order stays after voluntary withdrawal,
and what happens to an order and submitted action when a charge exhausts a budget. Orders should be
cancelled on budget exhaustion. Q's prompt and mechanism give different descriptions of when a round
ends, so its released implementation needs checking before a replication is called faithful
\cite{Q}.

P's energy has no natural exchange rate with Q's auction values. The peers proposed fixing the
tariff after a pilot and examining a declared range of \(\lambda\). Q uses closed API models whose
parameter counts are unavailable, so a common per-token charge with \(\alpha=0\) is the workable
first test, rather than P's size-weighted rule. A finite-budget arm and a matched exogenous-removal
control would help separate a changed choice to quote from the loss of active traders. The peers had
not verified that Q's released records contain reliable per-call reasoning-token counts \cite{Q,P}.

Q's association between reasoning language and urgency covers Gemini Large and is not causal
evidence for this proposed effect \cite{Q}. P's results about model size and its job economy do not
establish an auction effect \cite{P}. The limited outside-source check found auctions for task
allocation, communication opportunities, or related token costs, but the cited descriptions did not
establish this particular fixed-value test \cite{Agora,DALA,Bandwidth,TokenStudy}. That check cannot
support a broad novelty claim.

\section{Why the thread stopped}

The verifier accepted the connection and the peers' \(7.5\) calculation, including the zero-surplus
sixth trade. It left the thread at \textsc{Draft} because the behavioural claims were explicitly
untested. The smallest restart is to check Q's released scheduler and token records, set one
affordable tariff in a pilot, then compare matched auctions with zero and disclosed per-token
charges, no binding budget, and the same withdrawal rule. The comparison would record gross surplus,
token use, net welfare, crossings, withdrawals, and calls per trade. Only after that comparison
could a finite-budget arm test whether exhaustion changes the answer.

\bibliographystyle{plainurl}
\bibliography{references}

\end{document}
